\subsection{Interaction between pointing and registration}

Table~\ref{tab:four-arms} compares all four combinations. To isolate stand-alone interventions, we also mapped level-pointed echoes with the midpoint origin and level axes. Ranking the stand-alone gains changed the preferred first intervention in five of 111 cells relative to ranking registration against the additional gain of leveling after registration. In each, leveling led by 0.31--0.85 percentage points. At 0.36125 Hz and 0.376 m, leveling alone gained 24.01 points, recovering 1,421 positions from echoes both policies acquired and 54 from newly acquired echoes. With spoke-time registration already applied, leveling recovered those 54 echoes and lost two, yielding 0.85 points. The interaction was $-23.16$ percentage points because either intervention corrected much of the same mapping error. At 90-degree bearing in the largest registration-gain wave, leveling with the midpoint origin instead moved 692 acquired positions outside tolerance and recovered three new echoes.

\begin{table}[tb]
\caption{Usable positions out of 6,144 scheduled visits per wave. Columns cross antenna pointing with registration time.}\label{tab:four-arms}
\centering\small
\begin{tabular}{@{}lrrrr@{}}
\toprule
 & \multicolumn{2}{c}{Body-fixed} & \multicolumn{2}{c}{Level}\\
\cmidrule(lr){2-3}\cmidrule(l){4-5}
Wave (Hz, m) & Midpoint & Spoke & Midpoint & Spoke \\
\midrule
0.3500, 0.150 & 5,712 & 5,712 & 6,144 & 6,144 \\
0.3325, 0.432 & 878 & 2,338 & 2,091 & 2,392 \\
0.36125, 0.376 & 1,163 & 2,586 & 2,638 & 2,638 \\
\bottomrule
\end{tabular}
\end{table}
